\documentclass[11pt,a4paper]{article}
\usepackage[T1]{fontenc}
\usepackage[utf8]{inputenc}
\usepackage{lmodern}
\usepackage[margin=26mm]{geometry}
\usepackage{amsmath,amssymb,amsthm,booktabs,microtype}
\usepackage{listings}
\usepackage[hidelinks]{hyperref}
\hypersetup{pdftitle={A first counterexample to the Archer-Bourne recurrence for the avoidance chain (312,54321:132)},
  pdfauthor={Xiaojian Zeng},
  pdfsubject={Exact finite disproof in permutation pattern avoidance}}
\newtheorem{theorem}{Theorem}
\newtheorem{lemma}[theorem]{Lemma}
\newtheorem{proposition}[theorem]{Proposition}
\DeclareMathOperator{\Av}{Av}
\DeclareMathOperator{\LDS}{LDS}


\title{A first counterexample to the Archer--Bourne recurrence\\
for the avoidance chain $(312,54321:132)$}
\author{Xiaojian Zeng\\[3pt]
  \small PBC School of Finance, Tsinghua University, 100190\\
  \small\href{mailto:xiaojianzeng@foxmail.com}{xiaojianzeng@foxmail.com}}
\date{}
\begin{document}
\maketitle

\begin{abstract}
Let $a_n$ count permutations avoiding the classical patterns $312$ and
$54321$ whose composition squares avoid $132$. We disprove the recurrence
$a_n=a_{n-1}+a_{n-2}+a_{n-3}+a_{n-4}+n-1$ proposed by Archer and Bourne
for $n\geq6$. Its first failure is at $n=11$: the exact count is $1265$,
whereas the proposed right-hand side is $1266$.
An elementary direct-sum reduction reduces the discrepancy to the count of
permutations ending in $1$. A complete finite certificate checks all $4181$
classical candidates of length $11$, excludes $4172$ by explicit square-pattern
witnesses, and leaves nine. We give those nine permutations and their squares,
the counts through length $11$, and reproducible verification instructions.
No corrected recurrence for arbitrary length is asserted.
\end{abstract}

\section{Statement and context}

Write $S_n$ for the permutations of $[n]=\{1,\ldots,n\}$, with one empty
permutation in $S_0$. A permutation contains a \emph{classical pattern}
$\sigma\in S_k$ if some subsequence of $k$ entries has the same relative
order as $\sigma$; it avoids $\sigma$ otherwise. The square is composition,
$\pi^2(i)=\pi(\pi(i))$. Define
\[
 C_n=\{\pi\in S_n:\pi\text{ avoids }312,54321,
                  \ \pi^2\text{ avoids }132\},\qquad a_n=|C_n|.
\]
B\'ona and Smith~\cite{BonaSmith2019} study strong pattern avoidance,
where a permutation and its square avoid the same pattern. Archer and
Geary~\cite{ArcherGeary2024} introduce chain avoidance, allowing different
restrictions on successive powers. In this notation $C_n$ avoids the chain
$(312,54321:132)$.

Archer and Bourne~\cite[Section~5]{ArcherBourne2026} conjecture that
\begin{equation}\label{eq:conjecture}
 a_n=a_{n-1}+a_{n-2}+a_{n-3}+a_{n-4}+n-1\qquad(n\geq6).
\end{equation}
Their Theorem~4.2 enumerates the neighboring class with $4321$ in place of
$54321$. Its proof decomposes a permutation at its minimum and identifies
an involution suffix. We use that method with an exact finite endpoint
certificate to refute~\eqref{eq:conjecture}.

\begin{theorem}\label{thm:main}
The counts for $0\leq n\leq11$ are
\[
 (a_0,\ldots,a_{11})=(1,1,2,5,11,22,45,89,174,338,655,1265).
\]
Equation~\eqref{eq:conjecture} holds for $6\leq n\leq10$ and first fails at
$n=11$, where
\[
 1265\ne655+338+174+89+10=1266.
\]
There are exactly nine members of $C_{11}$ ending in $1$, listed with their
squares in Table~\ref{tab:endpoints}.
\end{theorem}

This is a computer-assisted finite disproof. The contribution is the first
failing coefficient and the exhaustive certificate; the direct-sum method is
credited to the preceding work. A bounded primary-literature search on
12 September 2026 located no equivalent disproof. This is not a claim of
priority or exhaustive coverage of unpublished work.

\section{Reduction to permutations ending in one}

For $u\in S_r$ and $v\in S_s$, the direct sum $u\oplus v$ is the
concatenation of $u$ with $v_1+r,\ldots,v_s+r$. Put
$\delta_j=(j,j-1,\ldots,1)$ and, for $n\geq1$, let
\[
 b_n=|\{\pi\in C_n:\pi(n)=1\}|.
\]
The direct-sum argument below is the version of
\cite[proof of Theorem~4.2]{ArcherBourne2026} needed here.

\begin{lemma}\label{lem:sum}
Every nonempty member of $C_n$ has a unique expression
\[
 \rho\oplus\delta_{d_1}\oplus\cdots\oplus\delta_{d_\ell},
 \qquad \rho\in C_k,\quad \rho(k)=1,\quad 1\leq d_i\leq4,
\]
where $k\geq1$ and $\ell\geq0$. Consequently
\begin{equation}\label{eq:endpoint}
 a_n=b_n+\sum_{j=1}^{\min(4,n-1)}a_{n-j}\qquad(n\geq1).
\end{equation}
\end{lemma}

\begin{proof}
If the minimum of a $312$-avoider $\pi$ is in position $k$, every entry
before it is smaller than every entry after it. Otherwise entries $x,1,y$
with $x>y$ form $312$. Thus the first $k$ values are precisely $[k]$, and
$\pi=\rho\oplus\tau$ uniquely, with $\rho(k)=1$.
Neither $312$ nor $54321$ can cross a direct-sum boundary.
Each summand is an invariant interval for the map, so
\[
 (\rho\oplus\tau)^2=\rho^2\oplus\tau^2;
\]
see also~\cite[Lemma~2.1]{ArcherGeary2024}.

If $\tau$ is nonempty, any inversion in $\tau^2$, together with any earlier
entry of $\rho^2$, would form $132$. Therefore $\tau^2$ is increasing,
and hence is the identity. Conversely, if $\rho^2$ avoids $132$ and
$\tau^2$ is the identity, their direct sum avoids $132$: a mixed occurrence
would require an inversion in the right summand.

A $312$-avoiding involution is layered, that is, a direct sum of decreasing
blocks. To see this, suppose its minimum is in position $j$. The same
minimum decomposition isolates $[j]$. Involutivity forces its first entry
to be $j$. An ascent among the entries between this $j$ and the terminal
$1$ would form $312$ with $j$. The entire first block is consequently
$\delta_j$. Induction applies to the suffix. Conversely, such layered
permutations are $312$-avoiding involutions. They avoid $54321$ precisely
when every decreasing block has size at most four. Their block boundaries
are ascents, so the blocks are unique.

For a member of $C_n$ not ending in $1$, delete the last decreasing block
of this nonempty layered suffix. Its length $j$ belongs to $\{1,2,3,4\}$,
and the remainder is a nonempty member of $C_{n-j}$. Conversely, appending
$\delta_j$ to such a remainder preserves every restriction and gives a
permutation not ending in $1$. The final decreasing run has length exactly
$j$, so these cases are disjoint. This proves~\eqref{eq:endpoint}. The bound
$j\leq n-1$ excludes the empty remainder, which would double-count an
endpoint permutation.
\end{proof}

In particular, for $n\geq5$,
\begin{equation}\label{eq:residual}
 a_n-a_{n-1}-a_{n-2}-a_{n-3}-a_{n-4}=b_n.
\end{equation}
Thus the conjecture at each $n\geq6$ is equivalent to $b_n=n-1$.
The reduction is valid for every length and does not assume that square
avoidance is inherited by arbitrary subpermutations.

\section{The finite endpoint certificate}

An endpoint permutation has the unique form
\[
 \pi=(\alpha_1+1,\ldots,\alpha_{n-1}+1,1),\qquad\alpha\in S_{n-1}.
\]
A terminal minimum cannot participate in $312$. A four-term descent in
$\alpha$ extends to a five-term descent in $\pi$; conversely any five-term
descent in $\pi$ supplies a four-term descent in $\alpha$.
The classical endpoint domain is therefore a copy of
$\Av_{n-1}(312,4321)$, before the square restriction is imposed.

\begin{lemma}\label{lem:domain}
Let $f_h(m)$ count $312$-avoiders of length $m$ with longest decreasing
subsequence of length at most $h$. For $h\geq0$, $f_h(0)=1$; for $m>0$,
$f_0(m)=0$. For $h,m\geq1$,
\begin{equation}\label{eq:domain}
 f_h(m)=\sum_{k=0}^{m-1}f_{h-1}(k)f_h(m-1-k).
\end{equation}
In particular $|\Av_{10}(312,4321)|=f_3(10)=4181$.
\end{lemma}

\begin{proof}
The minimum decomposition writes the permutation uniquely as
$(\alpha+1),1,(\beta+k+1)$, where $|\alpha|=k$ and
$|\beta|=m-1-k$. Both $\alpha$ and $\beta$ avoid $312$, and this condition
also suffices for the concatenation to avoid $312$.
No decreasing subsequence crosses the direct-sum boundary after $1$, so
the longest decreasing-subsequence length is
\[
 \max\{\LDS(\alpha)+1,\LDS(\beta)\},\qquad\LDS(\varnothing)=0.
\]
This proves the recurrence and its boundary cases. Integer evaluation gives
\[
 (f_3(0),\ldots,f_3(10))=(1,1,2,5,13,34,89,233,610,1597,4181).
\]
\end{proof}

\begin{proposition}\label{prop:finite}
The exact endpoint counts are
\[
 (b_1,\ldots,b_{11})=(1,1,2,3,3,5,6,7,8,9,9).
\]
\end{proposition}

\begin{proof}[Finite certificate proof]
For each $1\leq n\leq11$, the endpoint census visits every permutation
of $\{2,\ldots,n\}$ and appends $1$. It rejects a prefix if it contains
$312$ or $4321$, and tests every triple in the composition square of each
remaining endpoint. Thus it visits exactly $(n-1)!$ tuples, without
symmetry reduction or an assumed counting recurrence.
The direct index tests are given in Appendix~\ref{sec:reproduce}.

At length $11$ the certificate records $4181$ distinct endpoint permutations,
their squares, and either an increasing triple of one-based indices forming
$132$ in the square or a null entry. A separate checker imports no census
code. It checks bijectivity, the terminal $1$, distinctness, and original
$312$ and $54321$ avoidance, using the relative ranks of all relevant
subsequences. By Lemma~\ref{lem:domain}, these $4181$ verified distinct
members exhaust the classical domain.

The checker recomputes every square. It verifies all $4172$ exclusion triples
by testing $q_i<q_k<q_j$ for $i<j<k$, where $q=\pi^2$. For each of the nine
null entries it tests all triples and finds no $132$ occurrence. These nine
permutations are displayed in Table~\ref{tab:endpoints}.
Partitioning the certificate by first entry gives classical counts
$1597,1597,987$ at first entries $2,3,4$, respectively; the surviving counts
are $5,0,4$.
The exhaustive smaller endpoint lists give the remaining values in
Table~\ref{tab:counts}. All comparisons use exact integers.
\end{proof}

\begin{table}[tbp]
\centering\small
\caption{All members of $C_{11}$ ending in $1$, and their composition squares.
Each comma separates two values; $10$ and $11$ are single entries.}
\label{tab:endpoints}
\begin{tabular}{@{}ll@{}}
\toprule
$\pi$ & $\pi^2$ \\
\midrule
 $(2,3,4,5,6,7,8,9,10,11,1)$ & $(3,4,5,6,7,8,9,10,11,1,2)$ \\
 $(2,4,3,5,6,7,8,9,10,11,1)$ & $(4,5,3,6,7,8,9,10,11,1,2)$ \\
 $(2,5,4,3,6,7,8,9,10,11,1)$ & $(5,6,3,4,7,8,9,10,11,1,2)$ \\
 $(2,5,4,7,6,3,8,9,10,11,1)$ & $(5,6,7,8,3,4,9,10,11,1,2)$ \\
 $(2,5,4,7,6,9,8,3,10,11,1)$ & $(5,6,7,8,9,10,3,4,11,1,2)$ \\
 $(4,3,5,6,7,8,9,10,11,2,1)$ & $(6,5,7,8,9,10,11,2,1,3,4)$ \\
 $(4,3,5,7,6,9,8,11,10,2,1)$ & $(7,5,6,8,9,10,11,1,2,3,4)$ \\
 $(4,3,6,5,8,7,10,9,2,11,1)$ & $(5,6,7,8,9,10,11,2,3,1,4)$ \\
 $(4,3,6,7,5,8,9,10,11,2,1)$ & $(7,6,8,9,5,10,11,2,1,3,4)$ \\
\bottomrule
\end{tabular}
\end{table}

\begin{table}[tbp]
\centering\small
\caption{Exact counts through length $11$. The endpoint count $b_0$ is not defined.}
\label{tab:counts}
\begin{tabular}{@{}lrrrrrrrrrrrr@{}}
\toprule
$n$ & 0 & 1 & 2 & 3 & 4 & 5 & 6 & 7 & 8 & 9 & 10 & 11 \\
\midrule
$a_n$ & 1 & 1 & 2 & 5 & 11 & 22 & 45 & 89 & 174 & 338 & 655 & 1265 \\
$b_n$ & --- & 1 & 1 & 2 & 3 & 3 & 5 & 6 & 7 & 8 & 9 & 9 \\
\bottomrule
\end{tabular}
\end{table}

This separation of a generated witness from its checker follows the
certifying-algorithm framework of McConnell, Mehlhorn, N\"aher, and
Schweitzer~\cite{McConnellEtAl2011}. A related finite-certificate architecture
appears in Zeng, Liu, Ratnavelu, and Ong Seng Huat~\cite{ZengEtAl2026}:
their standard-library verifier reconstructs a finite domain and checks stored
coloring witnesses. Here the domain-completeness argument is
Lemma~\ref{lem:domain}, and the exclusion witnesses are triples of indices.
The cited geometric results are not assumptions in this permutation proof.

\begin{proof}[Proof of Theorem~\ref{thm:main}]
Starting with $a_0=1$, substitute Proposition~\ref{prop:finite} into
\eqref{eq:endpoint}. This gives the values of $a_n$ in Table~\ref{tab:counts},
including the six specified initial values. For $6\leq n\leq10$ the endpoint
counts are $b_n=n-1$, whereas $b_{11}=9\ne10$.
Equation~\eqref{eq:residual} proves that $11$ is the first failing order
in the conjecture's stated range.
\end{proof}

\section{Scope and computational assistance}

The failed identity is the universally quantified recurrence
\eqref{eq:conjecture}. The endpoint certificate determines only the listed
finite range; it gives no corrected formula for $b_n$ or $a_n$ at arbitrary
length. The neighboring $4321$ theorem is unaffected.

Besides the endpoint census and certificate checker, a separately written
program exhaustively visited every member of $S_n$ for $0\leq n\leq11$,
a total of $43\,954\,714$ permutations. It reproduced the counts, the full
$1265$-member set at length $11$, and the complete classical endpoint domain.
This provides an algorithmic cross-check without using the direct-sum
recurrence to generate the permutations.

OpenAI Codex assisted with literature searches, proof development, exact
enumeration and verification code, a separate mathematical review pass, and
manuscript preparation. The separate review pass used a fresh model context
and newly written verification code; it was not journal peer review or a
proof-assistant formalization. The mathematical evidence consists of the
arguments and exact finite checks described above. No floating-point fit,
random sampling, or language-model confidence is used in the proof.

\appendix
\section{Reproduction of the finite checks}\label{sec:reproduce}

The companion files preserve the endpoint census, every accepted endpoint
through length $11$, and the complete length-$11$ certificate. From a fresh
extraction of the companion archive, using Python~3 and its standard library,
run the following commands in the displayed order:
\begin{lstlisting}[basicstyle=\small\ttfamily,columns=fullflexible,keepspaces=true,showstringspaces=false,breaklines=true]
python3 audit/referee_verify.py
python3 evidence/endpoint_exhaustive.py
python3 evidence/verify_endpoint_certificate.py
\end{lstlisting}
The first command performs the full symmetric-group cross-check and imports
no endpoint or earlier enumeration implementation. The other two regenerate
the endpoint census and check its certificate. Use a working copy, since
these commands write JSON outputs. The census uses the following tests,
with zero-based array indices:
\begin{align*}
 312\text{ in }p&\Longleftrightarrow \exists i<j<k:\ p_i>p_k>p_j,\\
 4321\text{ in }p&\Longleftrightarrow \exists i<j<k<\ell:\ p_i>p_j>p_k>p_\ell,\\
 132\text{ in }q&\Longleftrightarrow \exists i<j<k:\ q_i<q_k<q_j,
 \qquad q_i=p_{p_i-1}.
\end{align*}
The prefix tests apply to the tuple from $\{2,\ldots,n\}$; the square
test applies after appending $1$. The factorial loop visits all tuples and
checks its iteration count against $(n-1)!$ for every length.

The file \path{evidence/endpoint_exhaustive.json} contains the per-length
counts and accepted lists. In \path{evidence/endpoint_classical_n11.json},
each record has fields \texttt{permutation}, \texttt{square}, and
\texttt{square\_132\_indices}; a null last field designates an accepted
square. The checker computes~\eqref{eq:domain} itself, verifies each record,
and reconstructs the counts by the proved reduction. Comparisons with the
additional archived enumerations occur after the endpoint computation and
are not used to generate it.

The companion README records input dependencies, output files, and the
reason for running the integrity-checking command before regeneration.

\bibliographystyle{plain}
\bibliography{references}
\end{document}
